\documentclass[11pt]{article}
\usepackage[margin=25mm]{geometry}
\usepackage{amsmath,amssymb,amsthm}
\usepackage[T1]{fontenc}
\usepackage[colorlinks=true,urlcolor=blue]{hyperref}
\newtheorem{theorem}{Theorem}
\setlength{\parindent}{0pt}\setlength{\parskip}{4pt}
\setlength{\emergencystretch}{3em}
\begin{document}
\begin{center}
{\Large A torsion-free path violates the alternating tree sum}\\[5pt]
{\large Conjecture 00000004228}\\[5pt]
zbw123456 \quad 14 October 2026
\end{center}
\textbf{Convention.} We use integral cellular homology and the usual
cellular spanning-tree counts: a zero-dimensional tree is a single vertex;
a one-dimensional spanning tree is a connected acyclic spanning subgraph.
For graphs, unweighted and torsion-weighted spanning-tree counts agree.
The original conjecture supplies no alternative definition of tree counts.

Let $X$ be the path $0\! -\! 1\! -\! 2$, with oriented edges $0\to1$ and
$1\to2$. Its complete cellular chain complex is
\[
0\longrightarrow C_1=\mathbb Z^2
\mathrel{\mathop{\longrightarrow}^{\partial}} C_0=\mathbb Z^3
\longrightarrow0,\qquad
\partial(a,b)=(-a,a-b,b).
\]
Write $s(x,y,z)=x+y+z$ for augmentation.

\begin{theorem}
All integral homology groups of $X$ are torsion-free, but
$\chi(X)\ne\tau_0(X)-\tau_1(X)$.
\end{theorem}
\begin{proof}
The equation $\partial(a,b)=0$ forces $a=b=0$, so $H_1(X)=0$.
Moreover, $s\partial=0$ and every $(x,y,z)$ with $x+y+z=0$ satisfies
\[
\partial(-x,z)=(x,y,z).
\]
Consequently $\operatorname{im}\partial=\ker s$ and $H_0(X)\cong\mathbb Z$.
All higher chain groups and homology groups are zero. This proves the
absence of torsion in every degree.

There are three vertices and two edges, so $\chi(X)=3-2=1$.
There are three zero-dimensional trees, one for each vertex.
A spanning subgraph is connected only if it contains both path edges,
and that subgraph is acyclic. Thus $\tau_0(X)=3$ and $\tau_1(X)=1$, giving
\[
\tau_0(X)-\tau_1(X)=3-1=2\ne1=\chi(X).
\]
This refutes the torsion-free implication; the unspecified correction
term in the presence of torsion is not needed.
\end{proof}

\textbf{Lean correspondence.} \texttt{EulerTrees.lean} defines the integral
chain groups and proves $\partial^2=0$, the kernel, and the image assertions.
The predicate \texttt{HomologyTorsionFree} states the exact quotient criterion
\[
q>0,\quad z\in Z_n,\quad qz\in B_n\quad\Longrightarrow\quad z\in B_n.
\]
Thus \texttt{homology\_torsion\_free} proves torsion-freeness directly, without
assuming homology ranks or omitting higher degrees.

For edge masks, \texttt{IsSpanningTree} requires integral connectivity
and a zero cycle kernel. \texttt{spanning\_tree\_iff} proves that only the full
mask is a tree. The mask list is exhaustive and duplicate-free; its filter
is proved equivalent to that predicate. The final theorem combines
all-degree torsion-freeness and the unequal Euler and tree counts.

\textbf{Scope and verification.} This is a counterexample under the stated
standard convention, not a result about an unspecified alternative formula.
Lean 4.19.0, core library only; the audit uses only \texttt{propext},
\texttt{Classical.choice}, and \texttt{Quot.sound}. Competition acceptance
has not been claimed.

\textbf{Source.} The official bilingual problem is included as
\texttt{conjecture.md}. For cellular trees and torsion weights, see
Duval, Klivans and Martin,
\href{https://jlmartin.ku.edu/papers/TreesChapter.pdf}{Simplicial and Cellular Trees}.
\end{document}
